\section{The Human Translation Stack and the Acoustic Toolkit}
\label{sec:humanmt}
\label{sec:foundationaudio}

Before asking what machines can make of a bark, we should be precise about
what machines already make of human speech and text. The last decade of
machine translation is not a single invention but a layered stack, and each
layer supplies a tool that animal decoding either inherits or lacks.
Figure~\ref{fig:genealogy} gives the timeline; Tables~\ref{tab:genealogy}
and~\ref{tab:audio} condense the two halves of the stack.

\begin{figure}[htbp]
\centering
\begin{tikzpicture}[x=1cm,y=1cm,
  box/.style={rectangle, rounded corners=3pt, draw=ink!40, fill=#1, text=ink,
    align=center, text width=2.2cm, minimum height=1.2cm, font=\scriptsize\bfseries},
  arr/.style={-{Stealth[length=2mm]}, thick, ink!60}
]
% milestone boxes (year + mechanism)
\node[box=paw!30]  at (1.1,1.1)  {2014\\seq2seq +\\attention};
\node[box=paw!30]  at (3.7,1.1)  {2017\\Transformer};
\node[box=sky!30]  at (6.3,1.1)  {2019--20\\BERT / GPT\\pretrain / LLM};
\node[box=sky!30]  at (8.9,1.1)  {2022\\NLLB\\multilingual};
\node[box=sage!30] at (11.5,1.1) {2023\\Whisper /\\Seamless};
\node[box=sage!30] at (14.1,1.1) {2023--24\\audio codecs\\+ audio-LLM};
% timeline axis
\draw[->, thick, ink!50] (-0.2,0) -- (15.4,0);
\foreach \x in {1.1,3.7,6.3,8.9,11.5,14.1}{
  \fill[ink!60] (\x,0) circle (1.6pt);
  \draw[arr] (\x,0.48) -- (\x,0.06);
}
% transfer band
\draw[arr, sagedeep] (7.6,-0.06) -- (7.6,-0.55);
\node[rectangle, rounded corners=4pt, draw=sagedeep, fill=sage!20, text=ink,
  align=center, text width=14.6cm, minimum height=0.9cm, font=\scriptsize\bfseries]
  at (7.6,-1.05) {transfer to animal decoding: AVES $\cdot$ NatureLM-audio $\cdot$ Perch $\cdot$ BEANS $\cdot$ CETI / DolphinGemma};
\end{tikzpicture}
\caption{\textbf{The human translation / Transformer genealogy, condensed.}
Each layer supplies a tool that animal decoding inherits (acoustic encoding,
discretization, sequence modeling, multilingual transfer) or lacks (a
reference corpus, shared grounding). The lower band marks where the toolkit
has already been pointed at animals.}
\label{fig:genealogy}
\end{figure}

Statistical MT treated translation as noisy-channel decoding over phrase
tables learned from parallel corpora such as Europarl~\cite{koehn2004europarl}
and OPUS~\cite{tiedemann2012opus}, toolkitted by Moses~\cite{koehn2005moses}.
The neural turn replaced hand-built phrases with end-to-end encoder--decoder
learning~\cite{sutskever2014seq2seq,cho2014rnn}, and the decisive ingredient
was attention, a soft learned alignment between source and target
positions~\cite{bahdanau2015attention,luong2015effective}. Convolutional
variants competed briefly~\cite{gehring2017conv}, and production systems
showed that deep recurrent attention models could be deployed at
scale~\cite{wu2016gnmt,ott2018scaling}. Two data-side inventions proved as
important as the architecture: subword segmentation (BPE), which composes
rare words from frequent pieces~\cite{sennrich2016bpe}, and back-translation,
which manufactures parallel signal from monolingual
text~\cite{sennrich2016improving,edunov2018backtrans}. Both speak directly to
the animal case: BPE is a recipe for discretizing a continuous signal into
reusable units, back-translation a recipe for learning when no parallel
corpus exists.

The Transformer dispensed with recurrence entirely~\cite{vaswani2017attention},
and its real legacy is that it made pretraining scalable: masked language
modeling (BERT)~\cite{devlin2019bert}, causal modeling
(GPT)~\cite{radford2018gpt,radford2019gpt2,brown2020gpt3}, and seq2seq
denoising (T5, BART)~\cite{raffel2020t5,lewis2020bart} turned massive
unlabeled corpora into general-purpose representations, with refinements
(RoBERTa, XLNet, ELECTRA, DeBERTa) establishing that the objective and data
regime matter more than architectural
cleverness~\cite{liu2019roberta,yang2019xlnet,clark2020electra,he2021deberta}.
Contextual embeddings had already shown the value of pretraining over
static vectors (word2vec, GloVe, ELMo,
ULMFiT)~\cite{mikolov2013word2vec,pennington2014glove,peters2018elmo,howard2018ulmfit}.
Scaling laws make returns predictable~\cite{kaplan2020scaling},
compute-optimal training clarifies the data/scale
trade-off~\cite{hoffmann2022chinchilla}, and the foundation-model paradigm
frames both opportunity and risk~\cite{bommasani2021foundation}. For
bioacoustics the lesson is that a large unlabeled corpus of vocalizations,
pretrained with a masked or contrastive objective, is the right starting
point.

Human MT's most encouraging result for animal decoding is that quality does
not require a bespoke parallel corpus per language pair. Shared-encoder
multilingual models enable zero-shot
transfer~\cite{johnson2017multilingual,firat2016multiway}; cross-lingual
pretraining (XLM, XLM-R, mBART, multilingual BERT) shares structure across
languages~\cite{conneau2019xlm,conneau2020xlmr,liu2020mbart,pires2019mbert};
mined parallelism (WikiMatrix, CCAligned) plus many-to-many training (M2M-100,
NLLB) brought hundreds of low-resource languages into
scope~\cite{schwenk2021wikimatrix,elkishky2020ccaligned,fan2021beyond,costajussa2022nllb,nllb2024scaling};
and fully unsupervised MT is possible from monolingual corpora
alone~\cite{artetxe2018unsupervised,lample2018unsupervised}. This is the
closest human analogue to the animal situation: a target with no parallel
corpus, where structure must be bootstrapped from raw signal and transferred
from better-resourced relatives.

Translation then moved end-to-end to speech: listen-and-translate proofs of
concept~\cite{berard2016listen}, direct speech
translation~\cite{weiss2017speech}, direct speech-to-speech
models~\cite{jia2019direct}, and unified multilingual multimodal systems
(SeamlessM4T, Seamless)~\cite{seamless2023m4t,seamless2024streaming}. The
acoustic-encoder-plus-semantic-decoder template is exactly what a genuine
animal translator would need, and it is engineerable when the source is
human language. What animals lack is the evaluation layer: human MT rests on
accumulated reference translations, BLEU~\cite{papineni2002bleu},
METEOR~\cite{denkowski2014meteor}, and learned metrics such as
COMET~\cite{rei2020comet}, validated against expert
judgment~\cite{freitag2021experts} and adjudicated annually at
WMT~\cite{bojar2018wmt,barrault2019wmt,kocmi2022wmt}. Every one of these
presupposes a ground-truth target sentence.

Open and proprietary LLMs (LLaMA, Llama~2, GPT-4, Gemini, Claude,
Gemma)~\cite{touvron2023llama,touvron2023llama2,openai2023gpt4,gemini2024,anthropic2024claude,gemma2024},
instruction tuning~\cite{ouyang2022instruct}, and chain-of-thought
reasoning~\cite{wei2022chain} supply the semantic and interactive layer;
vision-language models (CLIP, Flamingo, LLaVA,
BLIP-2)~\cite{radford2021clip,alayrac2022flamingo,liu2023llava,li2023blip2}
show how to bind a non-linguistic modality to language, the exact operation
a dog translator attempts when grounding a bark to a context.

A bark is not human speech, so the acoustic front-end matters. The decisive
idea is to learn a general representation from raw unlabeled waveform and
fine-tune on small labeled tasks: CPC~\cite{vandenoord2018cpc},
wav2vec~\cite{schneider2019wav2vec}, wav2vec~2.0~\cite{baevski2020wav2vec2},
HuBERT~\cite{hsu2021hubert}, WavLM~\cite{chen2022wavlm}, and
data2vec~\cite{baevski2022data2vec,baevski2023data2vec2}. Two extensions
matter directly for animals: wav2vec-U learns recognition without any
transcriptions~\cite{baevski2021wav2vecu}, and multilingual scaling (XLS-R,
MMS, USM) pushes one encoder across thousands of
languages~\cite{babu2022xlsr,pratap2023mms,zhang2023usm}. The
transcription-free recipe is the closest human precedent for a target with
no parallel corpus.

Recognition architectures solve the alignment problem in two regimes: CTC
marginalizes over all alignments between frames and
labels~\cite{graves2006ctc,graves2013speech,amodei2016deepspeech2}, while
attention-based listen-attend-spell learns a soft
alignment~\cite{chan2016las}; the Conformer marries convolution to
self-attention~\cite{gulati2020conformer}. For animal signals the problem is
the same in form but harder in substance, because there is no canonical
label sequence to align to. None of this would exist without shared data
and evaluation: LibriSpeech~\cite{panayotov2015librispeech}, Common
Voice~\cite{ardila2020commonvoice}, Multilingual
LibriSpeech~\cite{pratap2020mls}, TIMIT and
Switchboard~\cite{garofolo1993timit,godfrey1992switchboard}, and benchmarks
such as SUPERB~\cite{yang2021superb} and
CoVoST~\cite{wang2020covost,wang2021covost2}. The absence of a ``LibriSpeech
of barks'' is a concrete, fixable bottleneck.

A second lineage turns waveform into discrete tokens a language model can
predict: WaveNet~\cite{vandenoord2016wavenet},
VQ-VAE~\cite{vandenoord2017vqvae}, Jukebox~\cite{dhariwal2020jukebox},
SoundStream and EnCodec~\cite{zeghidour2021soundstream,defossez2022encodec},
AudioLM~\cite{borsos2023audiolm} and its music and sound
successors~\cite{agostinelli2023musiclm,copet2023musicgen,kreuk2023audiogen},
VALL-E~\cite{wang2023valle}, UniAudio~\cite{yang2024uniaudio}, and audio-LLMs
(AudioPaLM, SALMONN, Qwen-Audio)~\cite{rubenstein2023audiopalm,tang2024salmonn,chu2024qwenaudio}.
This is the discretization step a ``bark alphabet'' would require, the
strategy Project CETI applied to whale codas~\cite{sharma2024whale}.

The same representation revolution in vision, AlexNet and
ImageNet~\cite{krizhevsky2012alexnet,deng2009imagenet,russakovsky2015imagenet},
ResNet~\cite{he2016resnet}, and the Vision
Transformer~\cite{dosovitskiy2021vit}, is the lineage DogFACS-style facial
analysis rides on (Section~\ref{sec:biology}). And the toolkit has been
pointed at animals: BirdNET~\cite{kahl2021birdnet}, the reviews charting
computational bioacoustics~\cite{stowell2022review,stowell2014automatic},
AVES~\cite{hagiwara2023aves}, NatureLM-audio~\cite{robinson2024naturelm},
Perch~\cite{ghani2023perch}, BEANS~\cite{hagiwara2023beans}, and
OpenBEATs~\cite{openbeats2025}. The transfer is measured, not assumed:
fine-tuned animal-audio encoders report gains over from-scratch baselines on
BEANS and BirdSet. What none supplies is a reference translation; they
classify, embed, and retrieve, but do not render a bark as a sentence with
measured fidelity. The engineering of the front-end is largely solved; the
missing piece is semantic, not acoustic.

\begin{table}[htbp]
\centering
\caption{The human translation / Transformer genealogy and its transfer
value to animal decoding. ``Transfer value'' is our assessment.}
\label{tab:genealogy}
\footnotesize
\setlength{\tabcolsep}{4pt}
\begin{tabular}{@{}p{0.8cm}p{3.4cm}p{3.4cm}p{5.6cm}@{}}
\toprule
\textbf{Year} & \textbf{Milestone} & \textbf{Key mechanism} & \textbf{Transfer value to animals} \\
\midrule
2006 & CTC~\cite{graves2006ctc} & alignment-free labeling & label variable-length vocalizations \\
2014 & seq2seq~\cite{sutskever2014seq2seq,cho2014rnn} & encoder--decoder & signal sequence $\to$ symbol sequence \\
2015 & attention~\cite{bahdanau2015attention} & soft alignment & align bark segments to labels \\
2016 & BPE / back-translation~\cite{sennrich2016bpe,sennrich2016improving} & subwords; synthetic parallelism & discretize; bootstrap without parallel data \\
2017 & Transformer~\cite{vaswani2017attention} & self-attention & the universal backbone \\
2017 & multilingual NMT~\cite{johnson2017multilingual} & shared encoder, zero-shot & one encoder across species \\
2019--20 & BERT / GPT / T5~\cite{devlin2019bert,brown2020gpt3,raffel2020t5} & masked / causal pretraining & foundation encoder from raw signal \\
2021 & ViT / CLIP~\cite{dosovitskiy2021vit,radford2021clip} & cross-modal transformers & bind non-linguistic modality to language \\
2022 & NLLB / Chinchilla~\cite{costajussa2022nllb,hoffmann2022chinchilla} & mined parallel; compute-optimal & low-resource recipe; data-vs-scale trade-off \\
2023 & LLMs / Whisper / Seamless~\cite{openai2023gpt4,radford2023whisper,seamless2023m4t} & scale; weak-supervision ASR & semantic layer; robust acoustic encoder \\
2024 & NLLB (Nature)~\cite{nllb2024scaling} & 200 languages deployed & proof that low-resource scale works \\
\bottomrule
\end{tabular}
\end{table}

\begin{table}[htbp]
\centering
\caption{The speech/audio foundation-model toolkit and its transfer value.
SSL = self-supervised; codec = discrete acoustic tokens.}
\label{tab:audio}
\footnotesize
\setlength{\tabcolsep}{4pt}
\begin{tabular}{@{}p{4.6cm}p{1.6cm}p{7.0cm}@{}}
\toprule
\textbf{Model / family} & \textbf{Type} & \textbf{Transfer value to animals} \\
\midrule
wav2vec~2.0 / HuBERT / data2vec~\cite{baevski2020wav2vec2,hsu2021hubert,baevski2022data2vec} & SSL & foundation bark encoder from unlabeled audio \\
wav2vec-U~\cite{baevski2021wav2vecu} & SSL & learn without a parallel corpus \\
XLS-R / MMS / USM~\cite{babu2022xlsr,pratap2023mms,zhang2023usm} & SSL & one encoder across species \\
CTC / LAS / Conformer~\cite{graves2006ctc,chan2016las,gulati2020conformer} & ASR & map variable-length signal to labels \\
Whisper~\cite{radford2023whisper} & ASR & robust off-the-shelf acoustic encoder \\
LibriSpeech / Common Voice / SUPERB~\cite{panayotov2015librispeech,ardila2020commonvoice,yang2021superb} & data & template for a ``bark corpus'' \\
SoundStream / EnCodec~\cite{zeghidour2021soundstream,defossez2022encodec} & codec & discretize vocalizations into tokens \\
AudioLM / VALL-E / UniAudio~\cite{borsos2023audiolm,wang2023valle,yang2024uniaudio} & codec-LM & search for sequence structure \\
AudioPaLM / SALMONN / Qwen-Audio~\cite{rubenstein2023audiopalm,tang2024salmonn,chu2024qwenaudio} & audio-LLM & bind acoustic signal to language \\
AVES / NatureLM / Perch / BEANS~\cite{hagiwara2023aves,robinson2024naturelm,ghani2023perch,hagiwara2023beans} & bioacoustic & cross-species transfer, already realized \\
\bottomrule
\end{tabular}
\end{table}
